\section{Executive proof sketch on the public theorem spine}
\label{app:executive_proof_sketch}

\paragraph{Reserve-material status.}
This file is retained only as a routing aid. It adds no mathematics, carries no unique proof step, and is not part of the live proof burden of the hardened four-root suite. If any theorem-local implication appears only here, it must be moved back into one of the four compiled roots before the suite is considered synchronized.




This appendix follows the same ten public packets named in Section~\ref{sec:roadmap} and certified in Appendix~\ref{app:referee_audit}.
It is intentionally theorem-facing: each paragraph below names only the public theorem package used at that packet and suppresses side branches, historical scaffolding, and sample-sector audits.
The purpose is not to replace the detailed proofs, but to let a reader follow the theorem stack in the same order in which the paper actually uses it.
The packet summaries below are synchronized with the Packet~1 certification note, the Packet~3/8 transport ledger, the Packet~5--7 interface map, and the Packet~9/10 provenance ledgers recorded in Appendix~\ref{app:referee_audit}.
Appendix~\ref{app:referee_audit}, \S\ref{subsec:analytic_gap_packet}, is the short hostile-referee packet for the remaining analytic mass-gap attack line and should be read together with the Packet~8 and Packet~9 summaries below.

\subsection{Packet 1: compact-simple ultraviolet gate}

Appendix~\ref{app:A1_compact_simple} proves Theorem~\ref{thm:A1_compact_simple_internal} and records its public Route~1 consequence in Corollary~\ref{cor:A1_compact_simple_spine}.
Inside that packet, the low-Casimir branch is now an internal shifted-torus Laplace theorem in the parabolic weight window, while the high-Casimir branch is the real-analytic Peter--Weyl theorem from the torus-strip coefficient-extraction package.
The front-end certification note at Appendix~\ref{app:A1_compact_simple}, \S\ref{subsec:A1_packet1_certification}, freezes that theorem chain, separates the group-only constants from the $(G,\rho)$-indexed ultraviolet constants in Table~\ref{tab:A1_packet1_constant_ledger}, and records explicitly that the exceptional-group case is closed inside the same compact-simple regime theorems with no separate exceptional-family appendix.
This packet is both the ultraviolet gate and the public compact-simple group-scope statement on the public theorem spine once an auxiliary Wilson representation $\rho$ has been chosen, and all later explicit Route~1 constants are read from this packet together with the synchronized $(G,\rho)$-ledgers in Appendix~K.
Packet~10 later proves, first, the finite faithful-Wilson uniformity theorem needed by the action-family universality machinery and, second, that the resulting continuum local theory is independent of the faithful choice of $\rho$, while Corollary~\ref{cor:faithful_rho_gap_ledger} records that the explicit Route~1 certificate remains $\rho$-indexed in the present manuscript.

\subsection{Packet 2: one-shot entrance into the KP/Dobrushin basin}

Appendix~F packages the ultraviolet entrance theorem as Proposition~\ref{prop:route1_QE1}.
That statement says that the exact blocked Wilson specification enters the KP/Dobrushin basin at the explicit block scale $B(\beta;G,\rho)$ and that the physical entrance scale remains bounded along the tuned dyadic trajectory.
This is the only load-bearing ultraviolet-to-infrared entry statement on Route~1.

\subsection{Packet 3: full fixed-lattice OS gap}

Once Packet~2 is available, Appendix~F upgrades the entrance statement to the fixed-lattice gap package Proposition~\ref{prop:route1_QE2}.
The public output is a tuned-sequence lower bound for the full fixed-lattice OS gap, not merely for a preferred observable family.
This is the lattice-side quantity later transported to the continuum.

\subsection{Packet 4: Lane~C flowed state and bounded base algebra}

Section~\ref{sec:rg} culminates in Theorem~\ref{thm:C3U}, with Corollary~\ref{cor:flow_state} giving the executive flowed-state statement and Corollary~\ref{cor:C3U_Aplus_tuned} identifying the tuned bounded positive-time base state.
This packet is the common interface between Route~1 and Lane~A: Route~1 needs the bounded OS-positive sector, while Lane~A starts from the already-constructed flowed continuum state.

\subsection{Packet 5: Lane~A breadth and finite-truncation inverse control}

Appendix~\ref{app:normalization_crosscheck}, \S\ref{subsec:LaneA_packet_interface_map}, now freezes the Packet~5--7 interface in Table~\ref{tab:LaneA_packet_interface_map}.
At Packet~5, the public spine uses Theorem~\ref{thm:LaneA_associated_graded_identity} and Theorem~\ref{thm:LaneA_block_coefficient_theorem} for the basis-free exact-dimension breadth step, and then Theorem~\ref{thm:canonical_block_one_shell_increment}, Theorem~\ref{thm:finite_truncation_inverse_control}, and Corollary~\ref{cor:finite_truncation_SFTE} for the one-shell transport, finite-truncation inverse control, and intrinsic SFTE/remainder package.
This packet does not yet construct a capwise mixed state or perform the inductive-union sharp-local extension; sample scalar, tensor, and loop sectors remain in the paper only as audited illustrations and normalization cross-checks.

\subsection{Packet 6: Lane~A closure on the full sharp-local algebra}

At finite cap, Theorem~\ref{thm:finite_cap_flowed_sector_bridge} is the first mixed-correlator bridge from the compatible family of finite-family packaged bounded capwise states furnished by Corollary~\ref{cor:finite_cap_mixed_family_packaging} to the local inverse-flow limits; Table~\ref{tab:LaneA_packet_interface_map} records that no earlier Packet~5 theorem should be read as proving this closure step.
No universal mixed bounded capwise state is invoked there: in Theorem~\ref{thm:finite_cap_flowed_sector_bridge}, the notation $\omega_{\mathrm{cap,bnd}}^{(u_{\mathrm{ren}})}$ is only the local shorthand declared after the relevant finite family has been fixed.
Its part~(i) is the unique downstream theorem-facing closure node for the mixed-channel Cauchy clause.
The natural local citation point for that seam is \hyperref[par:finite_cap_mixed_channel_certification]{the local certification paragraph immediately following Theorem~\ref{thm:finite_cap_flowed_sector_bridge}(i)}, which records the single-slot chain Theorem~\ref{thm:ins_to_OS} $\Rightarrow$ Corollary~\ref{cor:ins_to_OS_bounded_word} $\Rightarrow$ Proposition~\ref{prop:finite_cap_bounded_factor_insertion}, the actual-summand chain Theorem~\ref{thm:finite_truncation_inverse_control} $\Rightarrow$ Lemma~\ref{lem:finite_cap_mixed_channel_reduction} $\Rightarrow$ Theorem~\ref{thm:finite_cap_flowed_sector_bridge}(i), the mixed bounded-family packaging/covariance node Corollary~\ref{cor:finite_cap_mixed_family_packaging}, the fixed-word bookkeeping through $K_W$, $C_W^{\mathrm{ins}}$, and the finite truncation data, and the displayed slotwise-to-functional $o(1)$ handoff in one place.
Appendix~\ref{app:referee_audit}, \S\ref{subsec:packet6_mixed_channel_note}, gives a synchronized theorem-location guide for the same seam.
Theorem~\ref{thm:LaneA_closure} is then the capwise-extension node with the finite-cap OS structure in part~(b), Corollary~\ref{cor:LaneA_bounded_state_restriction} is the bounded-state compatibility node consumed later on the QE3 seam, Corollary~\ref{cor:LaneA_full_scope} is the separate inductive-union passage, and Theorem~\ref{thm:LaneA_bounded_base_cyclic} is the first theorem that adds bounded-base cyclicity.
This is the only load-bearing Lane~A closure chain on the public theorem spine.

\subsection{Packet 7: bounded-base cyclicity}

The decisive continuity statement for the mass-gap transfer is Theorem~\ref{thm:LaneA_bounded_base_cyclic}: the original bounded positive-time base algebra remains cyclic in the full reconstructed sharp-local Hilbert space.
This is the packet that lets QE3 act on the full sharp-local vacuum sector rather than only on the bounded-domain closure.

\subsection{Packet 8: Wilson-to-continuum transport and the continuum vacuum gap}

On the public theorem spine, the continuum mass-gap route is compressed into one theorem-level package and recorded locally in \Cref{rem:route1_packet38_transport_certification}.
Theorem~\ref{thm:Wilson_patch_transport} closes the same-scale Wilson-versus-patched seam, Theorem~\ref{thm:continuum_tightness} identifies the Wilson bounded-state / dyadic OS limits, Theorem~\ref{thm:continuum_time_OS_upgrade} upgrades those limits to a continuum-time OS semigroup, Theorem~\ref{thm:phys_gap_noncollapse} converts the tuned full fixed-lattice gap into the full sharp-local Hilbert-space gap, and Corollary~\ref{cor:route1_QE3_gap} records the final continuum sharp-local vacuum gap.
That local ledger also records that Theorem~\ref{thm:weak_window_sufficiency} is the only live weak-window certificate on this route, that Corollary~\ref{cor:LaneA_bounded_state_restriction} is the first Lane~A theorem-level compatibility input actually consumed by the transport chain, that this corollary is not the inductive-union theorem, that Theorem~\ref{thm:LaneA_closure}(b) is the finite-cap sharp-local OS-structure input, that Corollary~\ref{cor:LaneA_full_scope} is the separate inductive-union passage, that Corollary~\ref{cor:sharp_base_dense_in_full_H} and Lemma~\ref{lem:QE3_core_topology} provide the later density/core handoff, and that Lane~B remains downstream only.
Appendix~\ref{app:referee_audit}, \S\ref{subsec:analytic_gap_packet}, compresses the corresponding constant route and the three classic uniformity seams into one theorem-facing packet.
This packet is the only load-bearing proof of \Cref{thm:main}(D).

\subsection{Packet 9: dense-core clustering, axioms, Wightman reconstruction, and non-triviality}

The dense-core Euclidean clustering reformulation is collected in Theorems~\ref{thm:route1_E2_internal} and~\ref{thm:R1_HS}; it is downstream from Packet~8 and not used to prove the gap.
Appendix~\ref{app:axioms_package} then records the local-field existence dossier in theorem form: Theorem~\ref{thm:OS_axioms_package} packages the Euclidean/OS properties from Theorem~\ref{thm:LaneA_closure}(b), Corollary~\ref{cor:LaneA_full_scope}, Theorem~\ref{thm:LaneA_temperedness_verification}, and Theorem~\ref{thm:LaneA_bounded_base_cyclic}, together with the standard OS reconstruction step on the full positive-time algebra; Theorem~\ref{thm:OS_to_Wightman} gives the Wightman/Haag--Kastler reconstruction; Corollary~\ref{cor:field_correspondence} identifies the reconstructed local fields with the renormalized gauge-invariant local differential polynomials represented by Lane~A; and Theorem~\ref{thm:nontriviality_package} consumes the separate witness export Theorem~\ref{thm:LaneA_nontriviality_witness}. The provenance ledger in Appendix~\ref{app:axioms_package}, \S\ref{subsec:packet9_provenance_ledger}, separates the imported OS reconstruction / analytic-continuation / uniqueness background from the earlier internal Euclidean input before the theorem packet is read linearly.
This claim-facing packet packages \Cref{thm:main}(B) and the theorem-facing existence claims in \Cref{thm:main}(C).

\subsection{Packet 10: exact local-net endpoint}

Appendix~\ref{app:global_form_resolution} fixes the final theorem-level local-net endpoint.
Theorem~\ref{thm:faithful_Wilson_uniform_hypotheses} first closes the finite faithful-Wilson comparison-family hypothesis gap needed by the action-family universality theorems.
Theorem~\ref{thm:faithful_Wilson_rep_universality} then proves that two faithful Wilson regularizations of the same compact simple group yield canonically isomorphic continuum sharp-local theories once the same admissible flow renormalization prescription is imposed; the only imported background at that step is the standard OS uniqueness clause used to upgrade equality of the Schwinger families to unitary equivalence of the OS/Wightman reconstructions.
Corollary~\ref{cor:faithful_rho_gap_ledger} records that the explicit Route~1 entrance/gap ledger remains indexed by the chosen faithful ultraviolet regularization.
Theorem~\ref{thm:global_form_local_endpoint} and Corollary~\ref{cor:global_form_main_scope} then identify the exact local-net endpoint and show that a nonfaithful ultraviolet choice merely descends to the effective quotient $G_{\mathrm{eff}}=G/\ker(\rho)$. The provenance ledger in Appendix~\ref{app:global_form_resolution}, \S\ref{subsec:packet10_provenance_ledger}, separates the earlier internal universality / descent input from the imported OS uniqueness clause used in faithful-Wilson universality, the new endpoint theorems, and the AQFT commentary remarks.
Packet~10 is therefore a theorem of local-net scope; the short closing remark in Appendix~\ref{app:global_form_resolution} records separately how this theorem-level endpoint matches the standard local-field reading.

\subsection{Route~1 summary in theorem-level terms}

Route~1 is the chain Theorem~\ref{thm:A1_compact_simple_internal} $\Longrightarrow$ Proposition~\ref{prop:route1_QE1} $\Longrightarrow$ Proposition~\ref{prop:route1_QE2} $\Longrightarrow$ Theorem~\ref{thm:Wilson_patch_transport}, Theorem~\ref{thm:continuum_tightness}, Theorem~\ref{thm:continuum_time_OS_upgrade}, Theorem~\ref{thm:phys_gap_noncollapse}, and Corollary~\ref{cor:route1_QE3_gap}.
The local certification point is \Cref{rem:route1_packet38_transport_certification}: it records Theorem~\ref{thm:weak_window_sufficiency} as the only live weak-window certificate on this chain, identifies Corollary~\ref{cor:LaneA_bounded_state_restriction} as the first Lane~A theorem-level compatibility input actually consumed by the transport route, states explicitly that this corollary is not the inductive-union theorem, records Theorem~\ref{thm:LaneA_closure}(b) as the finite-cap sharp-local OS-structure input, records Corollary~\ref{cor:LaneA_full_scope} as the separate inductive-union passage, points to Corollary~\ref{cor:sharp_base_dense_in_full_H} and Lemma~\ref{lem:QE3_core_topology} for the later density/core handoff, and states that Lane~B is not used anywhere on the proof of the gap itself.

\subsection{Lane~A summary in theorem-level terms}

Lane~A is the chain Theorem~\ref{thm:C3U} and Corollary~\ref{cor:C3U_Aplus_tuned} $\Longrightarrow$ Theorem~\ref{thm:LaneA_associated_graded_identity}, Theorem~\ref{thm:LaneA_block_coefficient_theorem}, Theorem~\ref{thm:canonical_block_one_shell_increment}, Theorem~\ref{thm:finite_truncation_inverse_control}, and Corollary~\ref{cor:finite_truncation_SFTE} $\Longrightarrow$ Theorem~\ref{thm:finite_cap_flowed_sector_bridge} $\Longrightarrow$ Theorem~\ref{thm:LaneA_closure}, Corollary~\ref{cor:LaneA_bounded_state_restriction}, Corollary~\ref{cor:LaneA_full_scope}, and Theorem~\ref{thm:LaneA_bounded_base_cyclic}, followed later by the Packet~9-facing exports Theorem~\ref{thm:LaneA_temperedness_verification} and Theorem~\ref{thm:LaneA_nontriviality_witness}.
Table~\ref{tab:LaneA_packet_interface_map} records the exact role split on that chain: Theorem~\ref{thm:finite_cap_flowed_sector_bridge} is the first mixed-correlator closure node, Theorem~\ref{thm:LaneA_closure}(b) is only the finite-cap OS-structure theorem, Corollary~\ref{cor:LaneA_bounded_state_restriction} is only the bounded-state compatibility node later consumed on the QE3 seam, Corollary~\ref{cor:LaneA_full_scope} is the separate inductive-union closure node, and Theorem~\ref{thm:LaneA_bounded_base_cyclic} is the first theorem that adds bounded-base cyclicity.
This is the only load-bearing route from the flowed probe algebra to the full sharp-local Hilbert space used by the continuum gap theorem and the Appendix~\ref{app:axioms_package} reconstruction dossier; the later Lane~A exports there are the separate temperedness and witness theorems just named.

\subsection{What lies outside the public spine}

Three collections of material now sit outside the public spine.
First, Lane~B is a downstream spectral-calculus reformulation fed by the already established gap.
Second, the classical-family appendices keep explicit worked one-plaquette audits for standard $\Sp/\Spin$ choices of $\rho$.
Third, the supplement retains archival Route~1 material moved off the live theorem chain during the hardening passes.
A referee should therefore treat Section~\ref{sec:roadmap} plus Appendices~\ref{app:referee_audit} and~\ref{app:executive_proof_sketch} as the first-pass certification packet, with \S\ref{subsec:analytic_gap_packet} of Appendix~\ref{app:referee_audit} serving as the short hostile-referee packet for the remaining analytic mass-gap attack line, and only then descend to the supporting layers.